\chapter{Building the MFJA Digital Twin}
\label{chap:digital-twin}

\section{Reconstruction strategy}

The work began with a minimal MFJA third-floor model, establishing a
recognisable, geometrically coherent environment before adding behaviour. The
third-floor mesh was imported and reoriented, the ground plane enlarged, and
Room~315 completed with doors and an interior window. Furniture, protective
guards, robot tables and laboratory assets followed.

Each asset required consistent units, origin, pose, visual material and
collision representation. CAD origins are often inconvenient manufacturing
datums, while Stereolithography (STL) files contain no unit metadata. Explicit
scale and orientation checks therefore aligned visual and collision geometry
with the intended room coordinates. Integration proceeded incrementally:

\begin{enumerate}
  \item load one asset in an isolated scene;
  \item verify scale and orientation against known room dimensions;
  \item establish a stable model frame;
  \item separate visual and collision geometry where necessary;
  \item place the asset in the full floor;
  \item add a static test for paths and SDF structure.
\end{enumerate}

The KUKA model followed this process, using a targeted, documented
collision-mesh configuration that preserved the required collision behaviour.

\section{World decomposition}

The final description package provides three principal world configurations:

\begin{itemize}
  \item \code{mfja\_3rd\_floor.world} for the shared third-floor environment;
  \item \code{room\_315\_only.world} for fast, focused rail, perception and
  language-integration work;
  \item isolated industrial-robot worlds assembled by launch for controller
  and geometry checks.
\end{itemize}

Because pose-update, entity-creation and removal service names include the
internal Gazebo world name, that name forms part of the interface. The bring-up
package explicitly passes both the selected world and its internal name, keeping
shuttle-controller services aligned with the active world; the runbook checks
the available services. \Cref{fig:room315-twin} shows the reconstruction.

\begin{figure}[p]
  \centering
  \captionsetup[subfigure]{skip=2pt}
  \begin{subfigure}[t]{0.96\textwidth}
    \centering
    \includegraphics[width=\linewidth]{early_full_floor.png}
    \caption{Earliest documented reconstruction stage produced during the internship.}
    \label{fig:early-floor-context}
  \end{subfigure}

  \vspace{1.5mm}
  \begin{subfigure}[t]{0.58\textwidth}
    \centering
    \includegraphics[width=\linewidth]{f01a_full_mfja_floor.png}
    \caption{Reconstructed third-floor context.}
    \label{fig:full-floor-context}
  \end{subfigure}

  \vspace{1.5mm}
  \begin{subfigure}[t]{0.75\textwidth}
    \centering
    \includegraphics[width=\linewidth]{f01b_room315_focus.png}
    \caption{Room~315 cells and robots.}
    \label{fig:room315-focus}
  \end{subfigure}
  \caption[Reconstruction progression from the floor model to Room~315.]
  {The reconstruction progressed from the earliest retained intermediate stage
  to the complete building context and the detailed Room~315 model.}
  \label{fig:room315-twin}
\end{figure}

\section{Repository refactoring}

As assets and launch components grew, the flat layout was refactored on 25
March into one repository containing several ROS packages. Its root ceased to
be a ROS package, and four packages received separate responsibilities:

\begin{description}
  \item[\code{mfja\_3rd\_floor\_description}] SDF, URDF, worlds, models,
  meshes and static description tests;
  \item[\code{mfja\_3rd\_floor\_bringup}] user-facing launch entry points and
  shared floor/Room~315 orchestration;
  \item[\code{mfja\_rail\_interfaces}] custom messages and the shuttle-add
  service;
  \item[\code{mfja\_robot\_control\_config}] robot configuration, rail
  controllers, planning, perception, learning tools, launches and tests.
\end{description}

This gives dependencies a clear direction: description assets remain separate
from planning code; interfaces remain independent of controller implementation;
bring-up composes the system without control algorithms; and control/configuration
depends on description and interfaces. It also supports selective builds with
\code{colcon}, the ROS~2 workspace build tool. \Cref{fig:packages} shows the
ownership and composition boundaries.

\begin{figure}[htbp]
  \centering
  \includegraphics[width=0.97\textwidth]{f09_ros2_package_architecture.png}
  \caption[ROS~2 package ownership and launch-composition dependencies.]
  {ROS~2 package ownership and composition. Description assets and
  typed rail contracts feed the control/configuration package; the bring-up
  package provides user-facing composition. Arrows represent dependency or
  launch composition, not runtime ROS message flow.}
  \label{fig:packages}
\end{figure}

\section{Model integration}

\subsection{Industrial robots}

The fixed-base robots are a KUKA KR6 R900 sixx, Stäubli TX2-60L, Yaskawa HC10
and Yaskawa HC10DT. Each model includes a link--joint kinematic tree, separate
visual and collision geometry, mounting geometry and a controller-facing joint
list. Integration preserves vendor-specific joint structure while exposing a
common higher-level command tool.

\subsection{TIAGo}

TIAGo has a mobile base, torso, arm and head, unlike the industrial arms; a
base-only variant supports focused tests. \Cref{chap:multi-robot} describes its
embodiment-specific command and feedback surfaces.

\subsection{Room and device assets}

Beyond the room shell, Room~315 has separate models for the static left and
right cell frames, rail paths, switches, shuttles, carried payloads, cameras and
visual markers. Markers aid inspection and can be hidden for clean camera
capture. The task uses two payload states: \emph{loaded} and \emph{empty}.

\section{Launch architecture}

The bring-up package exposes three stable entry points:
\path{full_floor.launch.py}, \path{room_315_only.launch.py} and
\path{single_industrial_robot.launch.py}. Shared Room~315 arguments are
declared once in \code{room\_315\_floor\_common.py}. This avoids configuration
drift in which the full-floor and room-only launch files gradually develop
different parameter names or defaults.

Launch choices cover graphical user interface (GUI) or headless mode (no
Gazebo graphical window), pause state, selected robots, kinematic shuttles,
Room~315 overhead cameras, visual markers, active shuttle identities, starting
slots and the initially loaded shuttle. Defaults leave the scene inactive or
readily inspectable; computationally expensive or research-specific nodes
require explicit activation. The launch layer also starts the ROS--Gazebo
bridges for entity creation, removal and pose updates.

\begin{lstlisting}[language=bash,caption={Focused Room~315 launch used as the basic transport test.},label={lst:room-launch}]
ros2 launch mfja_3rd_floor_bringup room_315_only.launch.py \
  robots:=none \
  start_paused:=false \
  gui:=true \
  enable_room315_kinematic_shuttles:=true
\end{lstlisting}

\section{Operational profiles}

The launch layer combines installed assets into repeatable operational
profiles: documented launch choices for a particular experiment. Operators can
isolate an integration boundary, control computational load and scale to the
shared environment without changing package ownership or public namespaces.
Profiles are documented in \path{runbook.html} and
\path{docs/QUICK_START_AND_FEATURE_GUIDE.md}; their verification status is
presented in \cref{chap:results}.

\section{Build and environment reproducibility}

The baseline is Ubuntu~24.04, ROS~2 Jazzy and Gazebo Harmonic in a standard
\code{colcon} workspace. Install rules cover the runtime scripts,
configuration, launch files, models and interfaces. Recorded campaigns and the
full-floor smoke test used executables and launch files from the built install
space. A smoke test is a narrow startup and interface check, not a performance
or reliability experiment. \Cref{app:reproducibility} specifies the environment
and build procedure.

\section{Outcome and model fidelity}

This phase produced usable full-floor and Room~315 simulation models,
selectable robots and a package structure supporting transport, perception,
language interpretation and planning. Recorded evaluations include one
all-robot full-floor cold start---new ROS and Gazebo processes rather than a
reused session---and the Room~315 integrated campaign. Later chapters report
transport, perception and result checking separately; \cref{chap:transport}
details the transport model and kinematic assumptions, and
\cref{chap:perspectives} discusses the next evaluation steps.
